\section{Two Books on One Desk}

Put the two books side by side and the first surprise is how much of the argument about them cannot be conducted from them. The 1962 \emph{Missale Romanum} is a single volume of about eleven hundred pages containing everything a priest needs at the altar: the general rubrics, the ceremonial directory, the calendar, the Order of Mass, the Canon, and every Mass formulary with its readings printed in place. The postconciliar \emph{Missale Romanum} in its third typical edition is a volume of comparable bulk containing almost the same things minus the readings, which now live in a separate book with its own introduction and its own three-year architecture. Neither book contains a word about which of them a given priest may lawfully use today; that is settled elsewhere, by acts that are not liturgical books at all. Neither contains an argument for itself. What each contains is a rite, printed with its rubrics, and the rubrics are extraordinarily specific about things the polemics rarely mention and silent about things the polemics treat as central.

This study reads the two books at their own loci. The method is pedestrian and the discipline is severe: for each point of comparison, find the place in each book where the matter is actually settled, quote or describe what stands there, and refuse to characterise either book from anything except its printed text and its own promulgating and governing acts. Where a claim commonly made about the two Missals does not survive that treatment, the claim is named and the text is shown. Where the claim survives, it is credited, including when it is uncomfortable.

Two consequences of the method should be stated at once, because they will disappoint different readers in different sections. The first is that this study cannot tell you which Missal is better. A book does not contain the measure of its own fruitfulness, and the question of what four decades of celebration have done to Catholic belief and practice is an empirical and historical question that no amount of textual collation can answer.\endnote{The wider question---what the Council mandated, how the reform was executed, and what the postconciliar statistical record will and will not support---is treated in a separate Triptych study, \emph{Council, Missal, and Crisis}, and is deliberately not reopened here. The present study is confined to the two books and the acts that govern them.} The second is that this study will not soften a difference to keep the peace. Several of the differences between the two books are large, and at least two of them are large in ways that the standard defence of the reform tends to pass over quickly. Saying so is not a verdict; it is the precondition of one.

\subsection{What the two objects are}

Precision about the objects is not pedantry here; almost every confused argument about ``the old Mass'' and ``the new Mass'' can be traced to an imprecision in naming one of them.

The older witness is the \emph{Missale Romanum, editio typica}, published at the Vatican Press in 1962. Its own opening decree, dated 23 June 1962 and signed by Cardinal Larraona as Prefect of the Sacred Congregation of Rites, explains why it exists: a new corpus of rubrics had been approved by John XXIII in the motu proprio \emph{Rubricarum instructum} and promulgated by the Congregation the following day, and so ``it could scarcely fail to happen'' that the same Congregation should prepare a new edition of the Roman Missal ``fully accommodated to the said code of rubrics.''\endnote{\emph{Missale Romanum}, editio typica (Vatican City, 1962), \emph{Decretum} of the Sacred Congregation of Rites, 23 June 1962, unnumbered front-matter leaf (artifact p.~2 of the registered facsimile): ``Novo rubricarum corpore, a Summo Pontifice Ioanne XXIII, Motu proprio \emph{Rubricarum instructum} diei 23 iulii anno 1960 approbato posteroque die a Sacra Rituum Congregatione promulgato, vix fieri non potuit quin eadem Sacra Rituum Congregatio novam Missalis romani editionem, ad dictum rubricarum codicem plane accommodatam, pararet.'' Collated visually at the page image, 25 July 2026. The Missal prints a discrepancy about the motu proprio's date: this decree gives 23 July 1960, while the Congregation's own general decree promulgating the code, printed a few leaves later and dated 26 July 1960, gives ``die 25 iulii huius anni.'' Only the latter is consistent with this decree's own ``postero die.'' The discrepancy is recorded in \texttt{research/edition-manifest.md} and left unresolved; no argument here depends on which date is right.} That sentence is worth holding on to. The 1962 book is not a survival that escaped reform; it is the product of a reform completed two years earlier, and it says so on its second leaf.

The newer witness is the \emph{Missale Romanum, editio typica tertia}, approved in 2000 and published in 2002, together with its 2008 emended reprint. It stands in a line whose earlier members---the 1969 \emph{Ordo Missae}, the 1970 first typical edition, the 1975 second---are distinct historical objects and are treated as such where they matter. The third edition is used here because it is the current normative Latin text, because it prints the 1969 apostolic constitution and the three promulgating decrees in its own front matter, and because comparing a current book with a current book removes one whole class of evasion from the argument.

Both books are Latin. Where this study quotes them, it quotes the Latin, because the Latin is the text that was promulgated and the only text in which the comparison can be made without importing a translator's decisions. The approved English of the postconciliar Missal and the text of the Lectionary are under copyright and are nowhere reproduced; they are cited by locus and described.\endnote{The rights position is set out in the terminal scope appendix. In brief: the Latin \emph{editiones typicae} are quoted; the ICEL English of the Roman Missal and the Lectionary pericopes are never reproduced at any length; public-domain English for the 1962 texts is available and is used where an English rendering is needed, always from an identified public-domain witness and never composed by this study.}

\subsection{The rule of the study}

Three working rules govern everything that follows, and they are stated here so that a reader can hold the study to them.

\textbf{First: a locus, or nothing.} Every substantive comparative claim is anchored to a place in a book---an article of the \emph{Institutio generalis}, a number in the \emph{Rubricae generales}, a rubric in the Order of Mass, an article of a motu proprio. A claim about ``the spirit'' of either book, or about what its compilers intended, or about what its use has produced, is either supported from a printed text or is not made.

\textbf{Second: keep the demonstrable and the contested apart.} A great deal about the two books is simply demonstrable: whether a prayer is present, what words it has, how many signs of the cross a rubric prescribes, which day outranks which. A great deal else is contested judgement: whether a shorter form says the same thing, whether repetition is redundancy or formation, whether a permission is a danger. This study separates them typographically and grammatically. Demonstrable matter is stated flatly. Contested judgement is marked as judgement, attributed where it belongs, and---in the one place where the study advances a judgement of its own---labelled as such and equipped with the finding that would overturn it.

\textbf{Third: take each side at its strongest.} The critical case against the postconciliar Missal is not the internet's version of it; it is the case made by a pope who legislated in its favour and by a constitution that the reform itself invokes. The case for the reform is not a shrug about pastoral necessity; it is an argument from the same phrase---\emph{ad pristinam sanctorum Patrum normam}---that Pius V used in 1570, deployed by the new Missal's own \emph{Institutio} with full awareness of what it was doing.\endnote{\emph{Institutio Generalis Missalis Romani} (2002), \emph{Proœmium} n.~6, in \emph{Missale Romanum}, editio typica tertia, artifact p.~10: the Council mandated that certain rites be restored ``ad pristinam sanctorum Patrum normam,'' ``iisdem videlicet usum verbis ac S.~Pius V in Constitutione Apostolica `Quo primum' inscriptis, qua anno 1570 Missale Tridentinum est promulgatum,'' and from that agreement of wording ``notari potest, qua ratione ambo Missalia romana, quamvis intercesserint quattuor sæcula, æqualem et parem complectantur traditionem.'' The article then adds the reform's own evaluative claim: ``intellegitur etiam, quam egregie ac feliciter prius perficiatur altero.''} A comparison that meets only the weaker version of either case has proved nothing about the stronger.
